% GENERATED FILE. Edit canonical lessons, then run npm run generate:books.
% canonical-source-hash: fnv1a64:71e2a80d18c539f8
% canonical-lessons: ML-C504-cavaru, ML-C504-valicceriyuka, ML-C504-kuminnukutuka, ML-C504-nasikkuka, ML-C504-punarupayogikkuka

\chapter{Litter}
\label{ch:ml-litter}
\hlchaptermodality{504}

\begin{chapteropening}
\textbf{By the end of this chapter:} \emph{I can say litter, to throw away, to pile up, to be ruined and to reuse.}

Complete the last lesson of chapter 504: I can say litter, to throw away, to pile up, to be ruined and to reuse.
\end{chapteropening}

\section[\texorpdfstring{\ml{ചവറ്} (cavaṟŭ)}{ചവറ് (cavaṟŭ)}]{\ml{ചവറ്} (cavaṟŭ) — litter; rubbish}

\label{lesson:ML-C504-cavaru}



\begin{quote}
Before the new one: say the Malayalam for to pollute, then the Malayalam for sewage.
\end{quote}

\subsection*{You'll want to know: \ml{ചവറ്}}
\textbf{\ml{ചവറ്}} — \emph{cavaṟŭ} — \textquotedblleft{}litter; rubbish\textquotedblright{}.

\begin{quote}
\ml{കടപ്പുറത്ത്} \ml{ധാരാളം} \ml{ചവറ്} \ml{ഉണ്ടായിരുന്നു}.
\end{quote}

\begin{quote}
\emph{kaṭappuṟattŭ dhārāḷaṁ cavaṟŭ uṇṭāyirunnu.}
\end{quote}

\begin{quote}
There was a lot of litter on the beach.
\end{quote}

Before some endings \ml{ചവറ്} doubles its last consonant: \ml{ചവറ്റിൽ}, in the rubbish.

\begin{cousinweb}
\textbf{\ml{ചവറ്}} (\emph{cavaṟŭ}) is native Dravidian, \textquotedblleft{}rubbish, sweepings\textquotedblright{}; DEDR also lists under it green leaves and rubbish used as manure (DEDR 2395).
\end{cousinweb}

\begin{grammarlens}[title={What you've built}]
One more word for getting around.
\end{grammarlens}

\subsection*{Your turn}
\begin{itemize}
\raggedright
  \item \emph{Say it:} \emph{cavaṟŭ}
  \item \emph{Say it:} \emph{cavaṟŭ}, once more
  \item \emph{Say it:} \emph{kaṭappuṟattŭ dhārāḷaṁ cavaṟŭ uṇṭāyirunnu.} — There was a lot of litter on the beach.
  \item \emph{Recall:} say the Malayalam for to pollute, then the Malayalam for sewage, then say \emph{cavaṟŭ} again
\end{itemize}

\begin{tcolorbox}[breakable,colback=teal!4,colframe=teal!35!black,title={Before you move on}]
What does \ml{ചവറ്} mean? (\textquotedblleft{}Litter; rubbish\textquotedblright{}.) Say it once more.
\end{tcolorbox}
\section[\texorpdfstring{\ml{വലിച്ചെറിയുക} (valicceṟiyuka)}{വലിച്ചെറിയുക (valicceṟiyuka)}]{\ml{വലിച്ചെറിയുക} (valicceṟiyuka) — to throw away; to fling}

\label{lesson:ML-C504-valicceriyuka}



\begin{quote}
Before the new one: say the Malayalam for sewage, then the Malayalam for litter.
\end{quote}

\subsection*{You'll want to know: \ml{വലിച്ചെറിയുക}}
\textbf{\ml{വലിച്ചെറിയുക}} — \emph{valicceṟiyuka} — \textquotedblleft{}to throw away; to fling\textquotedblright{}.

\begin{quote}
\ml{ആരോ} \ml{പ്ലാസ്റ്റിക്} \ml{കുപ്പി} \ml{പുഴയിലേക്ക്} \ml{വലിച്ചെറിഞ്ഞു}.
\end{quote}

\begin{quote}
\emph{ārō plāsṟṟikŭ kuppi puḻayilēkkŭ valicceṟiññu.}
\end{quote}

\begin{quote}
Someone flung a plastic bottle into the river.
\end{quote}

Past: \ml{വലിച്ചെറിഞ്ഞു}, \emph{valicceṟiññu}.

\ml{വലിച്ച്}, pulling, plus \ml{എറിയുക}, to throw: to throw away carelessly.

\begin{cousinweb}
\textbf{\ml{വലിച്ചെറിയുക}} (\emph{valicceṟiyuka}) is a native Dravidian compound of \ml{വലിക്കുക}, to pull (DEDR 5282), and \ml{എറിയുക}, to throw (DEDR 859).
\end{cousinweb}

\begin{grammarlens}[title={What you've built}]
One more word for getting around.
\end{grammarlens}

\subsection*{Your turn}
\begin{itemize}
\raggedright
  \item \emph{Say it:} \emph{valicceṟiyuka}
  \item \emph{Say it:} \emph{valicceṟiyuka}, once more
  \item \emph{Say it:} \emph{ārō plāsṟṟikŭ kuppi puḻayilēkkŭ valicceṟiññu.} — Someone flung a plastic bottle into the river.
  \item \emph{Recall:} say the Malayalam for sewage, then the Malayalam for litter, then say \emph{valicceṟiyuka} again
\end{itemize}

\begin{tcolorbox}[breakable,colback=teal!4,colframe=teal!35!black,title={Before you move on}]
What does \ml{വലിച്ചെറിയുക} mean? (\textquotedblleft{}To throw away; to fling\textquotedblright{}.) Say it once more.
\end{tcolorbox}
\section[\texorpdfstring{\ml{കുമിഞ്ഞുകൂടുക} (kumiññukūṭuka)}{കുമിഞ്ഞുകൂടുക (kumiññukūṭuka)}]{\ml{കുമിഞ്ഞുകൂടുക} (kumiññukūṭuka) — to pile up}

\label{lesson:ML-C504-kuminnukutuka}



\begin{quote}
Before the new one: say the Malayalam for litter, then the Malayalam for to throw away.
\end{quote}

\subsection*{You'll want to know: \ml{കുമിഞ്ഞുകൂടുക}}
\textbf{\ml{കുമിഞ്ഞുകൂടുക}} — \emph{kumiññukūṭuka} — \textquotedblleft{}to pile up\textquotedblright{}.

\begin{quote}
\ml{വഴിയോരത്ത്} \ml{മാലിന്യം} \ml{കുമിഞ്ഞുകൂടി}.
\end{quote}

\begin{quote}
\emph{vaḻiyōrattŭ mālinyaṁ kumiññukūṭi.}
\end{quote}

\begin{quote}
Rubbish piled up by the roadside.
\end{quote}

Past: \ml{കുമിഞ്ഞുകൂടി}, \emph{kumiññukūṭi}.

\ml{കുമിയുക}, to be heaped, plus \ml{കൂടുക}, to gather. For rubbish, work or money.

\begin{cousinweb}
\textbf{\ml{കുമിഞ്ഞുകൂടുക}} (\emph{kumiññukūṭuka}) is a native Dravidian compound of \ml{കുമിയുക}, to be heaped up (DEDR 1741), and \ml{കൂടുക}, to come together (DEDR 1882).
\end{cousinweb}

\begin{grammarlens}[title={What you've built}]
One more word for getting around.
\end{grammarlens}

\subsection*{Your turn}
\begin{itemize}
\raggedright
  \item \emph{Say it:} \emph{kumiññukūṭuka}
  \item \emph{Say it:} \emph{kumiññukūṭuka}, once more
  \item \emph{Say it:} \emph{vaḻiyōrattŭ mālinyaṁ kumiññukūṭi.} — Rubbish piled up by the roadside.
  \item \emph{Recall:} say the Malayalam for litter, then the Malayalam for to throw away, then say \emph{kumiññukūṭuka} again
\end{itemize}

\begin{tcolorbox}[breakable,colback=teal!4,colframe=teal!35!black,title={Before you move on}]
What does \ml{കുമിഞ്ഞുകൂടുക} mean? (\textquotedblleft{}To pile up\textquotedblright{}.) Say it once more.
\end{tcolorbox}
\section[\texorpdfstring{\ml{നശിക്കുക} (naśikkuka)}{നശിക്കുക (naśikkuka)}]{\ml{നശിക്കുക} (naśikkuka) — to be ruined; to perish}

\label{lesson:ML-C504-nasikkuka}



\begin{quote}
Before the new one: say the Malayalam for to throw away, then the Malayalam for to pile up.
\end{quote}

\subsection*{You'll want to know: \ml{നശിക്കുക}}
\textbf{\ml{നശിക്കുക}} — \emph{naśikkuka} — \textquotedblleft{}to be ruined; to perish\textquotedblright{}.

\begin{quote}
\ml{മഴയിൽ} \ml{വയലിലെ} \ml{നെല്ല്} \ml{നശിച്ചു}.
\end{quote}

\begin{quote}
\emph{maḻayil vayalile nellŭ naśiccu.}
\end{quote}

\begin{quote}
The paddy in the fields was ruined in the rain.
\end{quote}

Past: \ml{നശിച്ചു}, \emph{naśiccu}.

For crops, buildings and things left to rot.

\begin{cousinweb}
\textbf{\ml{നശിക്കുക}} (\emph{naśikkuka}) is built from the Sanskrit root \emph{naś}, \textquotedblleft{}to be lost, to perish\textquotedblright{} (CDIAL 7027), with the verb ending -\ml{ിക്കുക}.
\end{cousinweb}

\begin{grammarlens}[title={What you've built}]
One more word for getting around.
\end{grammarlens}

\subsection*{Your turn}
\begin{itemize}
\raggedright
  \item \emph{Say it:} \emph{naśikkuka}
  \item \emph{Say it:} \emph{naśikkuka}, once more
  \item \emph{Say it:} \emph{maḻayil vayalile nellŭ naśiccu.} — The paddy in the fields was ruined in the rain.
  \item \emph{Recall:} say the Malayalam for to throw away, then the Malayalam for to pile up, then say \emph{naśikkuka} again
\end{itemize}

\begin{tcolorbox}[breakable,colback=teal!4,colframe=teal!35!black,title={Before you move on}]
What does \ml{നശിക്കുക} mean? (\textquotedblleft{}To be ruined; to perish\textquotedblright{}.) Say it once more.
\end{tcolorbox}
\section[\texorpdfstring{\ml{പുനരുപയോഗിക്കുക} (punarupayōgikkuka)}{പുനരുപയോഗിക്കുക (punarupayōgikkuka)}]{\ml{പുനരുപയോഗിക്കുക} (punarupayōgikkuka) — to reuse; to recycle}

\label{lesson:ML-C504-punarupayogikkuka}



\begin{quote}
Before the new one: say the Malayalam for to pile up, then the Malayalam for to be ruined.
\end{quote}

\subsection*{You'll want to know: \ml{പുനരുപയോഗിക്കുക}}
\textbf{\ml{പുനരുപയോഗിക്കുക}} — \emph{punarupayōgikkuka} — \textquotedblleft{}to reuse; to recycle\textquotedblright{}.

\begin{quote}
\ml{ഞങ്ങൾ} \ml{പഴയ} \ml{കുപ്പികൾ} \ml{പുനരുപയോഗിക്കുന്നു}.
\end{quote}

\begin{quote}
\emph{ñaṅṅaḷ paḻaya kuppikaḷ punarupayōgikkunnu.}
\end{quote}

\begin{quote}
We reuse old bottles.
\end{quote}

Past: \ml{പുനരുപയോഗിച്ചു}, \emph{punarupayōgiccu}.

\ml{പുനർ}-, again, plus \ml{ഉപയോഗിക്കുക}, to use.

\begin{cousinweb}
\textbf{\ml{പുനരുപയോഗിക്കുക}} (\emph{punarupayōgikkuka}) is built from Sanskrit \emph{punar}, \textquotedblleft{}again\textquotedblright{}, and \emph{upayoga}, \textquotedblleft{}use\textquotedblright{} (Monier-Williams), with the verb ending.
\end{cousinweb}

\begin{grammarlens}[title={What you've built}]
One more word for getting around.
\end{grammarlens}

\subsection*{Your turn}
\begin{itemize}
\raggedright
  \item \emph{Say it:} \emph{punarupayōgikkuka}
  \item \emph{Say it:} \emph{punarupayōgikkuka}, once more
  \item \emph{Say it:} \emph{ñaṅṅaḷ paḻaya kuppikaḷ punarupayōgikkunnu.} — We reuse old bottles.
  \item \emph{Recall:} say the Malayalam for to pile up, then the Malayalam for to be ruined, then say \emph{punarupayōgikkuka} again
\end{itemize}

\begin{tcolorbox}[breakable,colback=teal!4,colframe=teal!35!black,title={Before you move on}]
What does \ml{പുനരുപയോഗിക്കുക} mean? (\textquotedblleft{}To reuse; to recycle\textquotedblright{}.) Say it once more.
\end{tcolorbox}
